% ============================================================
%  Beamer Presentation: Simulating Random Graph Models in Python
%  MA-300/MA-350/MA-360  Interim Presentation
%  Swansea University — Academic Year 2025/26
% ============================================================
\documentclass[t]{beamer}

\usepackage[english]{babel}
\usepackage{amsmath, amssymb, amsfonts}
\usepackage{booktabs}
\usepackage{array}
\usepackage{xcolor}

% ----- Theme ------------------------------------------------
\usetheme{CambridgeUS}
\usecolortheme{beaver}

% ----- Colour tweaks ----------------------------------------
\definecolor{swanseablue}{RGB}{0,62,116}
\definecolor{swanseagold}{RGB}{200,155,40}
\definecolor{lightgray}{RGB}{245,245,245}

\setbeamercolor{title}{fg=white, bg=swanseablue}
\setbeamercolor{frametitle}{fg=white, bg=swanseablue}
\setbeamercolor{structure}{fg=swanseablue}
\setbeamercolor{block title}{fg=white, bg=swanseablue}
\setbeamercolor{block body}{bg=lightgray}
\setbeamercolor{block title alerted}{fg=white, bg=darkred}
\setbeamercolor{block body alerted}{bg=lightgray}
\setbeamercolor{itemize item}{fg=swanseablue}
\setbeamercolor{itemize subitem}{fg=swanseagold}

% ----- Remove navigation symbols ---------------------------
\beamertemplatenavigationsymbolsempty

% ----- Metadata --------------------------------------------
\title[Random Graph Models]{Simulating Random Graph Models in Python}
\subtitle{Dissertation Presentation}
\author{Zachariah Markusson}
\institute[Swansea University]{Department of Mathematics\\
Swansea University}
\date{Academic Year 2025/26}

% ============================================================
\begin{document}

% ============================================================
%  SLIDE 1 — Title
% ============================================================
\begin{frame}
  \titlepage
\end{frame}

% ============================================================
%  SLIDE 2 — Motivation
% ============================================================
\section{Motivation}

\begin{frame}{Why Study Random Graph Models?}

\textbf{Networks are everywhere:}
\begin{itemize}
  \item Social networks, the internet, neural networks, citation graphs
  \item All share structural properties that cannot be captured by deterministic models
\end{itemize}

\medskip

\textbf{Key structural questions:}
\begin{itemize}
  \item How do large networks form and evolve?
  \item Why do real-world networks have short average path lengths (\emph{small-world property})?
  \item Why do a small number of nodes attract the majority of connections (\emph{hubs})?
\end{itemize}

\medskip

\begin{block}{Goal of this Project}
  Understand the \textbf{mathematical foundations} of three principal random graph models
  and explore their theoretical structural properties through rigorous analysis and simulation.
\end{block}

\end{frame}

% ============================================================
%  SLIDE 3 — Erdős–Rényi
% ============================================================
\section{Erd\H{o}s--R\'{e}nyi Model}

\begin{frame}{The Erd\H{o}s--R\'{e}nyi Model $G(n,p)$}

\begin{block}{Definition}
  Let $n \in \mathbb{N}$ be the number of vertices and $p \in [0,1]$.
  The random graph $G(n,p)$ is formed on vertex set $V = \{1,\ldots,n\}$
  by including each possible edge independently with probability $p$.
\end{block}

\medskip

\textbf{Fundamental properties:}
\begin{itemize}
  \item \textbf{Expected edges:}
        $\displaystyle\mathbb{E}[|E|] = p\binom{n}{2}$

  \item \textbf{Degree distribution:}
        $d(v) \sim \mathrm{Binomial}(n-1,\, p)$, so
        $\mathbb{E}[d(v)] = p(n-1)$

  \item \textbf{Connectivity threshold:}
        $p_c = \dfrac{\ln n}{n}$

        For $p > p_c$, the graph is \emph{almost surely} connected as $n \to \infty$;
        for $p < p_c$, isolated vertices persist almost surely.
\end{itemize}

\end{frame}

% ============================================================
%  SLIDE 4 — ER structural interpretation
% ============================================================

\begin{frame}{Erd\H{o}s--R\'{e}nyi: Structural Interpretation}
\footnotesize

\textbf{Phase transitions:}
\begin{itemize}
  \item Below $p_c$: the graph consists of many small, disconnected components.
  \item At $p = \frac{1}{n}$: a \emph{giant component} of size $\Theta(n)$ emerges.
  \item Above $p_c$: a single connected component dominates with high probability.
\end{itemize}

\smallskip

\textbf{Limitations as a real-world model:}
\begin{itemize}
  \item Degree distribution is \textbf{Binomial} (tightly concentrated around mean) — no heavy tails.
  \item \textbf{Low clustering:} the probability of two neighbours of $v$ being connected
        is simply $p$ — clustering coefficient $C \approx p \ll 1$ for sparse graphs.
  \item Does not replicate the high clustering observed in social or biological networks.
\end{itemize}

\begin{alertblock}{Key Takeaway}
  $G(n,p)$ gives sharp threshold theory but misses key real-network features:
  strong clustering and heavy-tailed degree distributions.
\end{alertblock}

\end{frame}

% ============================================================
%  SLIDE 5 — Watts–Strogatz
% ============================================================
\section{Watts--Strogatz Model}

\begin{frame}{The Watts--Strogatz Small-World Model}

\begin{block}{Construction}
  \begin{enumerate}
    \item Begin with a \textbf{regular ring lattice}: $n$ vertices arranged in a circle,
          each connected to its $k$ nearest neighbours ($k$ even, $k \ll n$).
    \item For each edge, \textbf{rewire} it with probability $\beta \in [0,1]$
          to a uniformly chosen random vertex (avoiding self-loops and duplicates).
  \end{enumerate}
\end{block}

\medskip

\textbf{Effect of the rewiring parameter $\beta$:}
\begin{itemize}
  \item $\beta = 0$: pure ring lattice — high clustering, large average path length.
  \item $\beta = 1$: effectively a random graph — low clustering, short path length.
  \item $0 < \beta \ll 1$: \textbf{small-world regime} — a few long-range shortcuts
        dramatically reduce average path length while clustering remains high.
\end{itemize}

\end{frame}

% ============================================================
%  SLIDE 6 — WS properties
% ============================================================

\begin{frame}{Watts--Strogatz: Clustering and Path Length}

\textbf{Clustering coefficient} (informally): the probability that two neighbours
of a vertex are themselves connected — a measure of local cohesion.

\medskip

\textbf{Why clustering stays high for small $\beta$:}
\begin{itemize}
  \item Most edges remain in place; the local lattice structure is preserved.
  \item The fraction of rewired edges is small, so triangles are not destroyed.
\end{itemize}

\medskip

\textbf{Why average path length drops sharply:}
\begin{itemize}
  \item Even a small number of random shortcuts act as \emph{bridges} between
        distant parts of the network.
  \item Reduces diameter from $O(n/k)$ on the lattice to $O(\log n)$.
\end{itemize}

\begin{alertblock}{The Small-World Phenomenon}
  High clustering (like a lattice) \emph{and} short path lengths (like a random graph)
  can coexist — precisely the regime observed in real-world networks.
\end{alertblock}

\end{frame}

% ============================================================
%  SLIDE 7 — Barabási–Albert
% ============================================================
\section{Barab\'{a}si--Albert Model}

\begin{frame}{The Barab\'{a}si--Albert Preferential Attachment Model}

\begin{block}{Construction}
  \begin{enumerate}
    \item Begin with a small initial connected graph on $m_0$ vertices.
    \item At each time step, add one new vertex with $m \leq m_0$ edges.
    \item Each new edge connects to an existing vertex $i$ with probability
          proportional to $i$'s current degree:
          \[
            P(i) \;=\; \frac{d_i}{\displaystyle\sum_{j} d_j}
          \]
  \end{enumerate}
\end{block}

\medskip

\textbf{Intuition:} \emph{``The rich get richer.''} \\
High-degree vertices attract more connections — a positive feedback mechanism.

\end{frame}

% ============================================================
%  SLIDE 8 — BA properties
% ============================================================

\begin{frame}{Barab\'{a}si--Albert: Scale-Free Networks}
\footnotesize

\textbf{Power-law degree distribution:}
\[
  P(k) \;\sim\; k^{-3}, \qquad k \to \infty
\]

This is derived rigorously via a mean-field or master-equation approach and holds
asymptotically as the network grows.

\smallskip

\textbf{What power-law implies structurally:}
\begin{itemize}
  \item \textbf{Heavy tail:} A small number of vertices (\emph{hubs}) accumulate
        a disproportionately large share of all edges.
  \item \textbf{Scale-free:} No characteristic degree scale exists —
        the distribution has no finite variance in the limit.
  \item \textbf{Robustness and vulnerability:} The network is highly robust
        to random failures (most nodes have low degree) but fragile under
        targeted attacks on hubs.
\end{itemize}

\begin{alertblock}{Key Takeaway}
  Preferential attachment explains heavy-tailed degrees and hubs seen in
  web and citation networks.
\end{alertblock}

\end{frame}

% ============================================================
%  SLIDE 9 — Comparison Table
% ============================================================
\section{Comparison}

\begin{frame}{Comparative Structural Properties}

\setlength{\tabcolsep}{4pt}
\renewcommand{\arraystretch}{1.15}
\begin{table}[h]
\scriptsize
\begin{tabular}{>{\bfseries}p{2.2cm} p{2.5cm} p{2.5cm} p{2.5cm}}
\toprule
 & \textbf{Erd\H{o}s--R\'{e}nyi} & \textbf{Watts--Strogatz} & \textbf{Barab\'{a}si--Albert} \\
\midrule
Mechanism      & Independent edge probability $p$ & Rewiring a ring lattice & Growth + preferential attachment \\
\addlinespace
Degree dist.   & Binomial$(n{-}1,p)$ & Approx.\ Poisson for small $\beta$ & Power law $P(k)\sim k^{-3}$ \\
\addlinespace
Clustering     & Low ($\approx p$) & High for small $\beta$ & Low–moderate \\
\addlinespace
Hubs           & None & None & Present (heavy tail) \\
\addlinespace
Typical uses   & Connectivity thresholds, percolation & Social, neural networks & Web, citation, scale-free nets \\
\bottomrule
\end{tabular}
\end{table}

\end{frame}

% ============================================================
%  SLIDE 10 — Implications, Limitations & Conclusion
% ============================================================
\section{Conclusion}

\begin{frame}{Theoretical Implications and Conclusions}
\footnotesize

\textbf{Theoretical implications:}
\begin{itemize}
  \item \textbf{ER:} Provides exact threshold phenomena via probabilistic methods,
        but is too homogeneous for most real networks.
  \item \textbf{WS:} Demonstrates that \emph{local structure} and \emph{global efficiency}
        are not mutually exclusive — a key insight for network science.
  \item \textbf{BA:} Shows that \emph{growth} and \emph{feedback} alone are sufficient
        to generate empirically observed heavy-tailed distributions.
\end{itemize}

\smallskip

\textbf{Limitations:}
\begin{itemize}
  \item WS does not produce power-law degree distributions.
  \item BA has low clustering — hybrid models (e.g., Holme--Kim) address this.
  \item All three assume \emph{undirected, unweighted} graphs.
\end{itemize}

\smallskip

\begin{block}{Summary}
  Each model captures a different structural mechanism.
  Together they motivate richer models for real-world complex networks.
\end{block}

\end{frame}

% ============================================================
\end{document}
